% 第26章 场论初步——习题解答
% 按 chapters/ch26.tex 中题目出现顺序编排。

\section*{26.1 散度和旋度——练习题}

\subsection*{练习题 1}
\begin{xsolution}
以下均假定所出现的一阶偏导数连续；关系式 (26.13)、(26.14) 还要求相应的二阶偏导数连续。

先证 (26.7)。设
\[
  f=f(x,y,z),\qquad g=g(x,y,z),
\]
则由偏导数的线性性，
\begin{align*}
  \nabla(\alpha f+\beta g)
  &\!=\vect{i}\frac{\partial(\alpha f+\beta g)}{\partial x}
   +\vect{j}\frac{\partial(\alpha f+\beta g)}{\partial y}
   +\vect{k}\frac{\partial(\alpha f+\beta g)}{\partial z}\\
  &=\alpha\nabla f+\beta\nabla g.
\end{align*}

设
\[
  \vect{a}=(a_1,a_2,a_3),\qquad \vect{b}=(b_1,b_2,b_3).
\]
则
\begin{align*}
  \nabla\cdot(\alpha\vect{a}+\beta\vect{b})
  &=\sum_{i=1}^3\frac{\partial(\alpha a_i+\beta b_i)}{\partial x_i}\\
  &=\alpha\nabla\cdot\vect{a}+\beta\nabla\cdot\vect{b},
\end{align*}
这就是 (26.8)。同理，由旋度三个分量都是一阶偏导数的线性组合，立即得到
\[
  \nabla\times(\alpha\vect{a}+\beta\vect{b})
  =\alpha\nabla\times\vect{a}+\beta\nabla\times\vect{b},
\]
即 (26.9)。

由乘积求导公式，
\begin{align*}
  \nabla(fg)
  &=\vect{i}(f_xg+fg_x)+\vect{j}(f_yg+fg_y)+\vect{k}(f_zg+fg_z)\\
  &=(\nabla f)g+f(\nabla g),
\end{align*}
故 (26.10) 成立。

再由
\begin{align*}
  \nabla\cdot(f\vect{a})
  &=\frac{\partial(fa_1)}{\partial x}
   +\frac{\partial(fa_2)}{\partial y}
   +\frac{\partial(fa_3)}{\partial z}\\
  &=f\left(\frac{\partial a_1}{\partial x}
   +\frac{\partial a_2}{\partial y}
   +\frac{\partial a_3}{\partial z}\right)
   +\left(f_xa_1+f_ya_2+f_za_3\right)\\
  &=f(\nabla\cdot\vect{a})+(\nabla f)\cdot\vect{a},
\end{align*}
得到 (26.11)。

对于 (26.12)，直接按分量计算。以第一分量为例，
\begin{align*}
  [\nabla\times(f\vect{a})]_1
  &=\frac{\partial(fa_3)}{\partial y}-\frac{\partial(fa_2)}{\partial z}\\
  &=f\left(\frac{\partial a_3}{\partial y}-\frac{\partial a_2}{\partial z}\right)
   +f_ya_3-f_za_2\\
  &=[f(\nabla\times\vect{a})+(\nabla f)\times\vect{a}]_1.
\end{align*}
其余两个分量完全同样，因此
\[
  \nabla\times(f\vect{a})
  =f(\nabla\times\vect{a})+(\nabla f)\times\vect{a}.
\]

最后，设 $\vect{a}=(a_1,a_2,a_3)$。由混合偏导数相等，
\begin{align*}
  \nabla\cdot(\nabla\times\vect{a})
  ={}&\frac{\partial}{\partial x}(a_{3,y}-a_{2,z})
  +\frac{\partial}{\partial y}(a_{1,z}-a_{3,x})
  +\frac{\partial}{\partial z}(a_{2,x}-a_{1,y})\\
  ={}&(a_{3,xy}-a_{3,yx})+(a_{2,zx}-a_{2,xz})+(a_{1,yz}-a_{1,zy})=0,
\end{align*}
故 (26.13) 成立；同理
\[
  \nabla\times(\nabla f)
  =(f_{zy}-f_{yz},\ f_{xz}-f_{zx},\ f_{yx}-f_{xy})=\vect{0},
\]
即 (26.14)。
\end{xsolution}

\subsection*{练习题 2}
\begin{xsolution}
以下按本节命题的通常语境，设向量场在 $\RR^3$ 上光滑；更一般地，只需定义域满足所用曲线、曲面均可在域内填充的相应拓扑条件。

先证命题 26.1.1(1)。若 $\vect{A}=\nabla u$，则对任意从 $P$ 到 $Q$ 的分段光滑曲线 $C$，由曲线积分基本定理，
\[
  \int_C\vect{A}\cdot\dd\vect{s}=u(Q)-u(P),
\]
所以积分与路径无关，$\vect{A}$ 是保守场；特别地，对任意闭曲线 $C$，
\[
  \oint_C\vect{A}\cdot\dd\vect{s}=0.
\]

反过来，若任意简单闭曲线上的环量为零，则任意两条从固定点 $P_0$ 到 $P$ 的分段光滑路径拼接后得到的闭曲线可分解为有限条简单闭曲线，从而两条路径上的积分相同。于是可定义
\[
  u(P)=\int_{P_0}^{P}\vect{A}\cdot\dd\vect{s},
\]
且定义与路径无关。取坐标方向的小线段可得
\[
  \frac{\partial u}{\partial x}=A_1,
  \qquad
  \frac{\partial u}{\partial y}=A_2,
  \qquad
  \frac{\partial u}{\partial z}=A_3,
\]
故 $\nabla u=\vect{A}$，即 $\vect{A}$ 为梯度场和保守场。

若 $\vect{A}=\nabla u$，由 (26.14) 有
\[
  \nabla\times\vect{A}=\nabla\times(\nabla u)=\vect{0},
\]
所以梯度场必为无旋场。反之若 $\nabla\times\vect{A}=\vect{0}$，则任意简单闭曲线 $C$ 在 $\RR^3$ 中可取一片分片光滑定向曲面 $S$ 以 $C$ 为边界。由 Stokes 公式，
\[
  \oint_C\vect{A}\cdot\dd\vect{s}
  =\iint_S(\nabla\times\vect{A})\cdot\vect{n}\,\dd S=0.
\]
结合上一段即得
\[
  \vect{A}\text{ 为梯度场}
  \Longleftrightarrow
  \vect{A}\text{ 为保守场}
  \Longleftrightarrow
  \vect{A}\text{ 为无旋场}
  \Longleftrightarrow
  \oint_C\vect{A}\cdot\dd\vect{s}=0.
\]

再证命题 26.1.1(2)。若 $\vect{A}=\nabla\times\vect{B}$，则由 (26.13)
\[
  \nabla\cdot\vect{A}
  =\nabla\cdot(\nabla\times\vect{B})=0,
\]
故旋度场必为无源场。

若 $\vect{A}$ 为 $\RR^3$ 上的无源场，则 $\RR^3$ 关于原点是星形区域，例题 26.1.2 保证存在光滑向量场 $\vect{B}$ 使
\[
  \vect{A}=\nabla\times\vect{B},
\]
故无源场也是旋度场。

另一方面，由 Gauss 公式，对任意分片光滑闭曲面 $\Sigma=\partial\Omega$，
\[
  \oiint_\Sigma\vect{A}\cdot\vect{n}\,\dd S
  =\iiint_\Omega\nabla\cdot\vect{A}\,\dd V.
\]
因此 $\nabla\cdot\vect{A}=0$ 时任意闭曲面的通量均为零。反之，若任意闭曲面的通量都为零，则对任意点 $p$ 及充分小的球 $B_\rho(p)$，
\[
  0=\oiint_{\partial B_\rho(p)}\vect{A}\cdot\vect{n}\,\dd S
  =\iiint_{B_\rho(p)}\nabla\cdot\vect{A}\,\dd V.
\]
若 $\nabla\cdot\vect{A}(p)>0$，由连续性可取足够小的 $\rho$ 使球内散度恒正，上式右端应为正，矛盾；散度为负的情形同理。因此
\[
  \nabla\cdot\vect{A}(p)=0.
\]
由于 $p$ 任意，$\vect{A}$ 为无源场。综上命题 26.1.1(2) 得证。
\end{xsolution}

\subsection*{练习题 3}
\begin{xsolution}
设 $V$ 沿一族含 $p_0$ 的分片光滑区域收缩到 $p_0$，即 $\operatorname{diam}V\to0$。

\begin{enumerate}[label=(\arabic*)]
  \item 由 Gauss 公式，
  \[
    \oiint_{\Sigma}\vect{A}\cdot\vect{n}\,\dd S
    =\iiint_V\operatorname{div}\vect{A}\,\dd V.
  \]
  因而
  \[
    \frac{1}{\abs{V}}\oiint_{\Sigma}\vect{A}\cdot\vect{n}\,\dd S
    =\frac{1}{\abs{V}}\iiint_V\operatorname{div}\vect{A}\,\dd V.
  \]
  由 $\operatorname{div}\vect{A}$ 在 $p_0$ 连续，右端是它在收缩区域上的积分平均值，故趋于 $\operatorname{div}\vect{A}(p_0)$。

  \item 写 $\vect{A}=(A_1,A_2,A_3)$。由 Gauss 公式逐分量应用，
  \begin{align*}
    \oiint_\Sigma(\vect{n}\times\vect{A})_1\,\dd S
    &=\oiint_\Sigma(n_2A_3-n_3A_2)\,\dd S\\
    &=\iiint_V\left(\frac{\partial A_3}{\partial y}-\frac{\partial A_2}{\partial z}\right)\,\dd V.
  \end{align*}
  另外两个分量同理，因此
  \[
    \oiint_\Sigma\vect{n}\times\vect{A}\,\dd S
    =\iiint_V\operatorname{curl}\vect{A}\,\dd V.
  \]
  除以 $\abs{V}$ 后再令 $\operatorname{diam}V\to0$，由旋度的连续性得到
  \[
    \lim_{\operatorname{diam}V\to0}
    \frac{1}{\abs{V}}\oiint_\Sigma\vect{n}\times\vect{A}\,\dd S
    =\operatorname{curl}\vect{A}(p_0).
  \]

  \item 同样逐分量使用 Gauss 公式：
  \[
    \oiint_\Sigma\varphi\vect{n}\,\dd S
    =\left(
      \iiint_V\varphi_x\,\dd V,
      \iiint_V\varphi_y\,\dd V,
      \iiint_V\varphi_z\,\dd V
    \right)
    =\iiint_V\operatorname{grad}\varphi\,\dd V.
  \]
  除以 $\abs{V}$ 并令区域收缩到 $p_0$，由 $\operatorname{grad}\varphi$ 的连续性即得
  \[
    \operatorname{grad}\varphi(p_0)
    =\lim_{\operatorname{diam}V\to0}\frac{1}{\abs{V}}
    \oiint_\Sigma\varphi\vect{n}\,\dd S.
  \]
\end{enumerate}
\end{xsolution}

\subsection*{练习题 4}
\begin{xsolution}
在 $r>0$ 时，
\[
  r=\sqrt{x^2+y^2+z^2},
  \qquad
  \nabla r=\frac{\vect{r}}{r}.
\]
因此由复合函数求导，
\[
  \operatorname{grad}f(r)=f'(r)\nabla r
  =\frac{f'(r)}{r}\vect{r}.
\]
又因为 $\nabla\cdot\vect{r}=3$，故由 (26.11)
\begin{align*}
  \operatorname{div}(f(r)\vect{r})
  &=f(r)\operatorname{div}\vect{r}+\nabla f(r)\cdot\vect{r}\\
  &=3f(r)+\frac{f'(r)}{r}\vect{r}\cdot\vect{r}\\
  &=3f(r)+rf'(r).
\end{align*}
最后，由 (26.12) 及 $\operatorname{curl}\vect{r}=0$，
\[
  \operatorname{curl}(f(r)\vect{r})
  =f(r)\operatorname{curl}\vect{r}+\nabla f(r)\times\vect{r}
  =\vect{0}.
\]
所以
\[
  \boxed{
  \operatorname{grad}f(r)=\frac{f'(r)}{r}\vect{r},\quad
  \operatorname{div}(f(r)\vect{r})=rf'(r)+3f(r),\quad
  \operatorname{curl}(f(r)\vect{r})=\vect{0}
  }.
\]
\end{xsolution}

\subsection*{练习题 5}
\begin{xsolution}
\begin{enumerate}[label=(\arabic*)]
  \item 由于
  \[
    \vect{r}=\nabla\left(\frac{x^2+y^2+z^2}{2}\right),
  \]
  由 (26.14) 有
  \[
    \operatorname{curl}\vect{r}=\vect{0}.
  \]

  \item 设 $\vect{c}=(c_1,c_2,c_3)$，则
  \[
    \vect{c}\times\vect{r}
    =(c_2z-c_3y,\ c_3x-c_1z,\ c_1y-c_2x).
  \]
  因而
  \begin{align*}
    \operatorname{curl}(\vect{c}\times\vect{r})
    &=(c_1-(-c_1),\ c_2-(-c_2),\ c_3-(-c_3))\\
    &=2\vect{c}.
  \end{align*}
\end{enumerate}
\end{xsolution}

\subsection*{练习题 6}
\begin{xsolution}
由练习题 4，条件
\[
  \operatorname{div}(f(r)\vect{r})=0
\]
等价于
\[
  rf'(r)+3f(r)=0,
  \qquad r>0.
\]
两边乘以 $r^2$，得到
\[
  r^3f'(r)+3r^2f(r)=0,
\]
即
\[
  \frac{\dd}{\dd r}\bigl(r^3f(r)\bigr)=0.
\]
因此
\[
  r^3f(r)=C,
\]
从而
\[
  \boxed{f(r)=\frac{C}{r^3}},\qquad r>0,
\]
其中 $C$ 为任意常数，$C=0$ 时即为零解。
\end{xsolution}

\subsection*{练习题 7}
\begin{xsolution}
由例题 26.1.3 的恒等式
\[
  \nabla\cdot(\vect{A}\times\vect{B})
  =\vect{B}\cdot(\nabla\times\vect{A})
  -\vect{A}\cdot(\nabla\times\vect{B}),
\]
而 $\vect{A},\vect{B}$ 均为无旋场，所以
\[
  \nabla\times\vect{A}=\nabla\times\vect{B}=\vect{0}.
\]
于是
\[
  \operatorname{div}(\vect{A}\times\vect{B})=0,
\]
即 $\vect{A}\times\vect{B}$ 是无源场。
\end{xsolution}

\section*{26.2 Laplace 算子与调和函数——练习题}

\subsection*{练习题 1}
\begin{xsolution}
\begin{enumerate}[label=(\arabic*)]
  \item 设 $f$ 二阶连续可微，则
  \[
    \nabla\times(\nabla f)
    =\left(
      f_{zy}-f_{yz},
      f_{xz}-f_{zx},
      f_{yx}-f_{xy}
    \right)=\vect{0}.
  \]

  \item 设 $\vect{a}=(a_1,a_2,a_3)$。考察第一分量：
  \begin{align*}
    [\nabla(\nabla\cdot\vect{a})]_1
    &=a_{1,xx}+a_{2,yx}+a_{3,zx},\\
    [\nabla\times(\nabla\times\vect{a})]_1
    &=\frac{\partial}{\partial y}(a_{2,x}-a_{1,y})
      -\frac{\partial}{\partial z}(a_{1,z}-a_{3,x})\\
    &=a_{2,xy}-a_{1,yy}-a_{1,zz}+a_{3,xz}.
  \end{align*}
  由混合偏导数相等，二者相减得到
  \[
    a_{1,xx}+a_{1,yy}+a_{1,zz}=\Delta a_1.
  \]
  对第二、第三分量作同样计算，便有
  \[
    \nabla(\nabla\cdot\vect{a})-\nabla\times(\nabla\times\vect{a})
    =(\Delta a_1,\Delta a_2,\Delta a_3)
    =\Delta\vect{a}.
  \]
\end{enumerate}
\end{xsolution}

\subsection*{练习题 2（第二 Green 恒等式）}
\begin{xsolution}
对 $u,v$ 使用第一 Green 恒等式，
\[
  \iiint_\Omega v\Delta u\,\dd V
  +\iiint_\Omega\nabla u\cdot\nabla v\,\dd V
  =\oiint_\Sigma v\frac{\partial u}{\partial n}\,\dd S.
\]
交换 $u,v$，得到
\[
  \iiint_\Omega u\Delta v\,\dd V
  +\iiint_\Omega\nabla v\cdot\nabla u\,\dd V
  =\oiint_\Sigma u\frac{\partial v}{\partial n}\,\dd S.
\]
第一式减去第二式，中间的梯度内积项抵消，故
\[
  \boxed{
  \iiint_\Omega(v\Delta u-u\Delta v)\,\dd V
  =\oiint_\Sigma\left(
  v\frac{\partial u}{\partial n}
  -u\frac{\partial v}{\partial n}
  \right)\,\dd S
  }.
\]
这就是第二 Green 恒等式。
\end{xsolution}

\subsection*{练习题 3}
\begin{xsolution}
在第一 Green 恒等式中令 $v=u$，得到
\[
  \iiint_\Omega u\Delta u\,\dd V
  +\iiint_\Omega\abs{\nabla u}^2\,\dd V
  =\oiint_\Sigma u\frac{\partial u}{\partial n}\,\dd S.
\]
因 $u$ 在 $\Omega$ 上调和，$\Delta u=0$，于是
\[
  \boxed{
  \iiint_\Omega\abs{\nabla u}^2\,\dd V
  =\oiint_\Sigma u\frac{\partial u}{\partial n}\,\dd S
  }.
\]

下面证明惟一性。设 $u_1,u_2$ 都在 $\Omega$ 内调和并在边界 $\Sigma$ 上取相同的值，令
\[
  w=u_1-u_2.
\]
则
\[
  \Delta w=0\quad\text{于 }\Omega,
  \qquad
  w=0\quad\text{于 }\Sigma.
\]
把上面的等式用于 $w$，有
\[
  \iiint_\Omega\abs{\nabla w}^2\,\dd V
  =\oiint_\Sigma w\frac{\partial w}{\partial n}\,\dd S=0.
\]
被积函数非负且连续，所以 $\nabla w=\vect{0}$ 于 $\Omega$。若 $\Omega$ 连通，则 $w$ 为常数；又 $w$ 在边界为零，故 $w\equiv0$。因此
\[
  u_1\equiv u_2,
\]
即调和函数在内部的值由边界值惟一确定。
\end{xsolution}

\subsection*{练习题 4}
\begin{xsolution}
由调和函数的球面平均值公式，对每个 $0<\rho\le R$，
\[
  \oiint_{\partial B_\rho(M_0)}u\,\dd S
  =4\pi\rho^2u(M_0).
\]
按以 $M_0$ 为中心的球坐标逐层积分，
\begin{align*}
  \iiint_{B_R(M_0)}u(x,y,z)\,\dd V
  &=\int_0^R\left(
    \oiint_{\partial B_\rho(M_0)}u\,\dd S
    \right)\dd\rho\\
  &=4\pi u(M_0)\int_0^R\rho^2\,\dd\rho\\
  &=\frac{4}{3}\pi R^3u(M_0).
\end{align*}
所以
\[
  \boxed{
  u(M_0)=\frac{1}{\frac43\pi R^3}
  \iiint_{B_R(M_0)}u(x,y,z)\,\dd x\,\dd y\,\dd z
  }.
\]
\end{xsolution}

\subsection*{练习题 5}
\begin{xsolution}
设 $u$ 在半径为 $R$ 的圆盘内调和。定义单位圆盘上的函数
\[
  v(\rho\cos\theta,\rho\sin\theta)
  =u(R\rho\cos\theta,R\rho\sin\theta),
  \qquad 0\le\rho<1.
\]
因为
\[
  \Delta v=R^2(\Delta u)\circ(R\,\cdot)=0,
\]
故 $v$ 在单位圆盘上调和。由命题 26.2.1，若 $s=r/R$，则
\begin{align*}
  u(r\cos\theta,r\sin\theta)
  &=v(s\cos\theta,s\sin\theta)\\
  &=\frac{1}{2\pi}\int_0^{2\pi}
    v(\cos\varphi,\sin\varphi)
    \frac{1-s^2}{1-2s\cos(\theta-\varphi)+s^2}\,\dd\varphi\\
  &=\frac{1}{2\pi}\int_0^{2\pi}
    u(R\cos\varphi,R\sin\varphi)
    \frac{R^2-r^2}{R^2-2Rr\cos(\theta-\varphi)+r^2}\,\dd\varphi.
\end{align*}
这正是 Poisson 积分公式 (26.20)。
\end{xsolution}

\subsection*{练习题 6}
\begin{xsolution}
记
\[
  q=\frac{r}{R},\qquad 0\le q<1.
\]
Poisson 核可展开为
\begin{align*}
  \frac{R^2-r^2}{R^2-2Rr\cos(\theta-\varphi)+r^2}
  &=\frac{1-q^2}{1-2q\cos(\theta-\varphi)+q^2}\\
  &=1+2\sum_{n=1}^\infty q^n\cos n(\theta-\varphi).
\end{align*}
对任意固定的 $0<\rho<R$，当 $r\le\rho$ 时，不但上式级数一致收敛，而且把它对 $x,y$ 求至二阶偏导所得级数也一致收敛，因为相应项可由常数倍的
\[
  n^2\left(\frac{\rho}{R}\right)^{n-2}
\]
控制，而该数项级数收敛。因此可以在积分号与级数号下逐项求二阶偏导。

对固定的 $\varphi$，每一项
\[
  r^n\cos n(\theta-\varphi)
\]
都是 $x,y$ 的调和多项式的线性组合，故其 Laplace 算子为零，常数项也显然调和。于是 Poisson 公式定义的函数满足
\[
  \Delta u(x,y)
  =\frac{1}{2\pi}\int_0^{2\pi}
  f(\varphi)\,\Delta_{x,y}
  \left[
  \frac{R^2-r^2}{R^2-2Rr\cos(\theta-\varphi)+r^2}
  \right]\dd\varphi=0
\]
于任意闭圆盘 $r\le\rho<R$。由于 $\rho$ 任意，$u$ 在整个开圆盘 $r<R$ 上调和。
\end{xsolution}

\subsection*{练习题 7}
\begin{xsolution}
给出一个同时适用于二维和三维的证明。设 $u$ 在区域 $\Omega$ 上调和，取任意 $x_0\in\Omega$，再取 $\varepsilon>0$ 使
\[
  \overline{B_{2\varepsilon}(x_0)}\subset\Omega.
\]
取一个非负、径向对称的 $C^\infty$ 函数 $\eta_\varepsilon$，其支集包含于 $B_\varepsilon(0)$ 且
\[
  \int\eta_\varepsilon(y)\,\dd y=1.
\]
对 $x\in B_\varepsilon(x_0)$，由调和函数的球面平均值公式，把积分按半径分层可得
\begin{align*}
  (u*\eta_\varepsilon)(x)
  &=\int_{B_\varepsilon(0)}u(x-y)\eta_\varepsilon(y)\,\dd y\\
  &=u(x)\int_{B_\varepsilon(0)}\eta_\varepsilon(y)\,\dd y
  =u(x).
\end{align*}
这里第二个等号的理由是：在每个固定半径的球面上，$u(x-y)$ 的平均值都等于 $u(x)$，而 $\eta_\varepsilon$ 在该球面上为常数。

卷积 $u*\eta_\varepsilon$ 是 $C^\infty$ 函数，因此 $u$ 在 $B_\varepsilon(x_0)$ 上无限次可微。由于 $x_0$ 任意，$u\in C^\infty(\Omega)$。
\end{xsolution}

\subsection*{练习题 8}
\begin{xsolution}
设定义域为连通开集 $G$。由 $f$ 与 $g\circ f$ 都调和，
\[
  \Delta f=0,
  \qquad
  \Delta(g(f))=0.
\]
利用复合函数求导公式，
\begin{align*}
  \Delta(g(f))
  &=g''(f)\abs{\nabla f}^2+g'(f)\Delta f\\
  &=g''(f)\abs{\nabla f}^2.
\end{align*}
故
\begin{equation}
  g''(f(x))\abs{\nabla f(x)}^2=0,
  \qquad x\in G.
  \tag{*}
\end{equation}

下面证明 $g''$ 在 $f(G)$ 上恒为零。反设存在 $x_0\in G$ 使
\[
  g''(f(x_0))\ne0.
\]
由 $g''$ 连续，可取含 $f(x_0)$ 的开区间 $I$，使 $g''$ 在 $I$ 上处处不为零。令 $U$ 为开集
\[
  \{x\in G\mid f(x)\in I\}
\]
中含 $x_0$ 的连通分支。由 (*)，在 $U$ 上必有 $\nabla f=0$，所以 $f$ 在 $U$ 上为常数，且等于 $f(x_0)$。

若 $U$ 在 $G$ 中有边界点 $p$，由连续性 $f(p)=f(x_0)\in I$。于是 $p$ 的某个小邻域仍包含于 $f^{-1}(I)$，且与 $U$ 相交，这与 $U$ 是 $f^{-1}(I)$ 的连通分支矛盾。因此 $U$ 在 $G$ 中既开又闭。由于 $G$ 连通，只能有 $U=G$，从而 $f$ 为常值函数，与题设矛盾。

因此 $g''(t)=0$ 对每个 $t\in f(G)$ 都成立。又因为连通集 $G$ 的连续像 $f(G)$ 是区间，所以在该区间上
\[
  g(t)=at+b.
\]
即 $g$ 在 $f(G)$ 上是线性函数（仿射函数）。
\end{xsolution}

\subsection*{练习题 9}
\begin{xsolution}
若问题有解 $u$，则由 Gauss 公式，
\begin{align*}
  \iiint_D f\,\dd x\,\dd y\,\dd z
  &=\iiint_D\Delta u\,\dd x\,\dd y\,\dd z\\
  &=\oiint_{\partial D}\nabla u\cdot\vect{n}\,\dd S\\
  &=\oiint_{\partial D}\frac{\partial u}{\partial n}\,\dd S\\
  &=\oiint_{\partial D}g\,\dd S.
\end{align*}
所以
\[
  \boxed{
  \iiint_D f\,\dd x\,\dd y\,\dd z
  =\oiint_{\partial D}g\,\dd S
  }
\]
是该 Neumann 问题有解的必要条件。
\end{xsolution}

\section*{26.3 参考题——第一组参考题}

\subsection*{第一组参考题 1}
\begin{xsolution}
记 $\varepsilon_{ijk}$ 为 Levi--Civita 符号，$\delta_{ij}$ 为 Kronecker 符号，并采用重复指标求和约定。用恒等式
\[
  \varepsilon_{ijk}\varepsilon_{klm}
  =\delta_{il}\delta_{jm}-\delta_{im}\delta_{jl}.
\]

\begin{enumerate}[label=(\arabic*)]
  \item 第 $i$ 个分量有
  \begin{align*}
    [\vect{A}\times(\nabla\times\vect{B})]_i
    &=\varepsilon_{ijk}A_j\varepsilon_{klm}\partial_lB_m\\
    &=A_m\partial_iB_m-A_l\partial_lB_i,
  \end{align*}
  同理
  \[
    [\vect{B}\times(\nabla\times\vect{A})]_i
    =B_m\partial_iA_m-B_l\partial_lA_i.
  \]
  再加上
  \[
    [(\vect{B}\cdot\nabla)\vect{A}]_i=B_l\partial_lA_i,
    \qquad
    [(\vect{A}\cdot\nabla)\vect{B}]_i=A_l\partial_lB_i,
  \]
  后两组项恰好抵消，留下
  \[
    A_m\partial_iB_m+B_m\partial_iA_m
    =\partial_i(\vect{A}\cdot\vect{B}).
  \]
  三个分量都成立，故
  \[
    \nabla(\vect{A}\cdot\vect{B})
    =\vect{A}\times(\nabla\times\vect{B})
    +\vect{B}\times(\nabla\times\vect{A})
    +(\vect{B}\cdot\nabla)\vect{A}
    +(\vect{A}\cdot\nabla)\vect{B}.
  \]

  \item 第 $i$ 个分量为
  \begin{align*}
    [\nabla\times(\vect{A}\times\vect{B})]_i
    &=\varepsilon_{ijk}\partial_j(\varepsilon_{klm}A_lB_m)\\
    &=\partial_m(A_iB_m)-\partial_l(A_lB_i)\\
    &=B_m\partial_mA_i-A_l\partial_lB_i
      +(\partial_mB_m)A_i-(\partial_lA_l)B_i.
  \end{align*}
  因此
  \[
    \nabla\times(\vect{A}\times\vect{B})
    =(\vect{B}\cdot\nabla)\vect{A}
    -(\vect{A}\cdot\nabla)\vect{B}
    +(\nabla\cdot\vect{B})\vect{A}
    -(\nabla\cdot\vect{A})\vect{B}.
  \]
\end{enumerate}
\end{xsolution}

\subsection*{第一组参考题 2}
% SOURCE-CHECK: lower/page-0396.jpg 原书题面为“F 为 G 上的光滑无源场”；chapters/ch26.tex 当前误录为“光滑无旋场”。以下按原书题意解答。
\begin{xsolution}
按原书题意，$\vect{F}$ 是无源场，即
\[
  \nabla\cdot\vect{F}=0.
\]
令 $\vect{r}=(x,y,z)$。对固定的 $t\in[0,1]$，在上一题 (2) 中取
\[
  \vect{X}(\vect{r})=t\vect{F}(t\vect{r}),
  \qquad
  \vect{Y}(\vect{r})=\vect{r}.
\]
因为
\[
  (\vect{X}\cdot\nabla)\vect{r}=\vect{X},
  \qquad
  \nabla\cdot\vect{r}=3,
\]
且
\[
  \nabla\cdot\vect{X}
  =t^2(\nabla\cdot\vect{F})(t\vect{r})=0,
\]
故上一题 (2) 给出
\begin{align*}
  \nabla\times[t\vect{F}(t\vect{r})\times\vect{r}]
  &=(\vect{r}\cdot\nabla)[t\vect{F}(t\vect{r})]
    -t\vect{F}(t\vect{r})+3t\vect{F}(t\vect{r})\\
  &=t^2(\vect{r}\cdot\nabla)\vect{F}(t\vect{r})
    +2t\vect{F}(t\vect{r})\\
  &=\frac{\dd}{\dd t}\left[t^2\vect{F}(t\vect{r})\right].
\end{align*}
于是可以在积分号下取旋度：
\begin{align*}
  \nabla\times\vect{A}(\vect{r})
  &=\int_0^1
    \nabla\times[t\vect{F}(t\vect{r})\times\vect{r}]\,\dd t\\
  &=\int_0^1\frac{\dd}{\dd t}
    [t^2\vect{F}(t\vect{r})]\,\dd t\\
  &=\vect{F}(\vect{r})-\lim_{t\to0^+}t^2\vect{F}(t\vect{r})\\
  &=\vect{F}(\vect{r}),
\end{align*}
因为 $\vect{F}$ 在原点附近光滑有界。因此
\[
  \boxed{\nabla\times\vect{A}=\vect{F}}.
\]
\end{xsolution}

\subsection*{第一组参考题 3}
\begin{xsolution}
题设给出
\[
  \nabla\times\vect{B}=\frac1r(\nabla r\times\vect{A}),
  \qquad r=\sqrt{x^2+y^2+z^2}.
\]
对两边取散度。由旋度的散度为零，
\[
  0=\nabla\cdot\left[\frac1r(\nabla r\times\vect{A})\right].
\]
利用乘积公式，
\begin{align*}
  0
  &=\nabla\left(\frac1r\right)\cdot(\nabla r\times\vect{A})
    +\frac1r\nabla\cdot(\nabla r\times\vect{A}).
\end{align*}
其中 $\nabla(1/r)$ 与 $\nabla r$ 平行，故第一项为零；又由
\[
  \nabla\cdot(\vect{P}\times\vect{Q})
  =\vect{Q}\cdot(\nabla\times\vect{P})
  -\vect{P}\cdot(\nabla\times\vect{Q}),
\]
取 $\vect{P}=\nabla r$、$\vect{Q}=\vect{A}$，并用
\[
  \nabla\times(\nabla r)=\vect{0},
\]
得到
\[
  \nabla\cdot(\nabla r\times\vect{A})
  =-\nabla r\cdot(\nabla\times\vect{A}).
\]
因此
\[
  \nabla r\cdot(\nabla\times\vect{A})=0.
\]

设 $L$ 位于半径为 $R$ 的球面上，取该球面上由 $L$ 围成的一片定向曲面 $S$。在 $S$ 上单位法向量可取
\[
  \vect{n}=\nabla r.
\]
于是由 Stokes 公式，
\[
  \oint_L\vect{A}\cdot\vect{\tau}\,\dd s
  =\iint_S(\nabla\times\vect{A})\cdot\vect{n}\,\dd S
  =0.
\]
故
\[
  \boxed{\oint_L\vect{A}\cdot\vect{\tau}\,\dd s=0}.
\]
\end{xsolution}

\subsection*{第一组参考题 4}
\begin{xsolution}
令
\[
  \vect{N}=\frac{\nabla F}{\abs{\nabla F}}.
\]
由二维散度定理（Green 公式），
\[
  \iint_D\nabla\cdot\vect{N}\,\dd x\,\dd y
  =\oint_C\vect{N}\cdot\vect{n}\,\dd s,
\]
其中 $\vect{n}$ 是 $D$ 的单位外法向量。

因为 $D=\{F>0\}$，在边界 $C$ 上，$\nabla F$ 指向 $F$ 增大的方向，即指向 $D$ 的内部。因此
\[
  \vect{n}=-\frac{\nabla F}{\abs{\nabla F}}=-\vect{N}.
\]
故
\[
  \vect{N}\cdot\vect{n}=-1.
\]
于是
\[
  \boxed{
  \iint_D\nabla\cdot\left(\frac{\nabla F}{\abs{\nabla F}}\right)\,\dd x\,\dd y
  =-\oint_C\dd s=-l
  }.
\]
\end{xsolution}

\subsection*{第一组参考题 5}
\begin{xsolution}
记 $M=(x_0,y_0,z_0)$，令 $\vect{\omega}$ 为单位球面上的单位向量。由球面面积的缩放关系，
\[
  F(R)=\frac{1}{4\pi}\oiint_{\partial B_1(0)}u(M+R\vect{\omega})\,\dd S_1.
\]
由 $u$ 在 $M$ 连续，$u(M+R\vect{\omega})\to u(M)$ 对 $\vect{\omega}$ 一致成立，因此
\[
  \lim_{R\to0^+}F(R)=u(M).
\]

下面求主要部分。由于 $u$ 在 $M$ 有连续二阶偏导数，Taylor 展开为
\[
  u(M+R\vect{\omega})
  =u(M)+R\nabla u(M)\cdot\vect{\omega}
  +\frac{R^2}{2}\vect{\omega}^{\mathsf T}H_u(M)\vect{\omega}
  +\littleo(R^2),
\]
其中余项对单位球面上的 $\vect{\omega}$ 一致。球面对称性给出
\[
  \frac1{4\pi}\oiint_{\partial B_1}\omega_i\,\dd S=0,
  \qquad
  \frac1{4\pi}\oiint_{\partial B_1}\omega_i\omega_j\,\dd S
  =\begin{cases}
  \frac13,&i=j,\\
  0,&i\ne j.
  \end{cases}
\]
因此
\begin{align*}
  F(R)-u(M)
  &=\frac{R^2}{2}\cdot\frac13
  \left(u_{xx}(M)+u_{yy}(M)+u_{zz}(M)\right)
  +\littleo(R^2)\\
  &=\frac{\Delta u(M)}{6}R^2+\littleo(R^2).
\end{align*}
若 $\Delta u(M)\ne0$，则无穷小量 $F(R)-u(M)$ 的主要部分为
\[
  \boxed{\frac{\Delta u(M)}6R^2}.
\]
\end{xsolution}

\subsection*{第一组参考题 6}
\begin{xsolution}
令
\[
  w=v-u.
\]
由于 $u=v$ 于 $\partial\Omega$，有 $w=0$ 于 $\partial\Omega$。展开
\begin{align*}
  \iiint_\Omega\abs{\nabla v}^2\,\dd V
  &=\iiint_\Omega\abs{\nabla u+\nabla w}^2\,\dd V\\
  &=\iiint_\Omega\abs{\nabla u}^2\,\dd V
    +2\iiint_\Omega\nabla u\cdot\nabla w\,\dd V
    +\iiint_\Omega\abs{\nabla w}^2\,\dd V.
\end{align*}
由第一 Green 恒等式，
\[
  \iiint_\Omega\nabla u\cdot\nabla w\,\dd V
  =\oiint_{\partial\Omega}w\frac{\partial u}{\partial n}\,\dd S
  -\iiint_\Omega w\Delta u\,\dd V=0,
\]
因为 $w=0$ 于边界且 $u$ 调和。故
\[
  \iiint_\Omega\abs{\nabla v}^2\,\dd V
  =\iiint_\Omega\abs{\nabla u}^2\,\dd V
  +\iiint_\Omega\abs{\nabla w}^2\,\dd V
  \ge\iiint_\Omega\abs{\nabla u}^2\,\dd V.
\]
即得所证不等式。
\end{xsolution}

\subsection*{第一组参考题 7}
\begin{xsolution}
令
\[
  x=r\cos\theta,
  \qquad
  y=r\sin\theta,
  \qquad
  0\le r<1.
\]
则
\[
  x\frac{\partial u}{\partial x}
  +y\frac{\partial u}{\partial y}
  =r\frac{\partial u}{\partial r}.
\]
因此所求积分为
\begin{equation}
  I=\int_0^1\int_0^{2\pi}r^2u_r(r,\theta)\,\dd\theta\,\dd r.
  \tag{1}
\end{equation}

固定 $0<r<1$，在圆盘 $B_r=\{x^2+y^2<r^2\}$ 上应用 Green 公式：
\begin{align*}
  \oint_{\partial B_r}
  \left(u_x\,\dd y-u_y\,\dd x\right)
  &=\iint_{B_r}\Delta u\,\dd x\,\dd y\\
  &=\iint_{B_r}\ee^{-(x^2+y^2)}\,\dd x\,\dd y\\
  &=2\pi\int_0^r\rho\ee^{-\rho^2}\,\dd\rho\\
  &=\pi(1-\ee^{-r^2}).
\end{align*}
在圆周 $\partial B_r$ 上有
\[
  \dd x=-r\sin\theta\,\dd\theta,
  \qquad
  \dd y=r\cos\theta\,\dd\theta,
\]
故左端等于
\[
  \int_0^{2\pi}r
  \left(u_x\cos\theta+u_y\sin\theta\right)\,\dd\theta
  =\int_0^{2\pi}r u_r(r,\theta)\,\dd\theta.
\]
于是
\[
  \int_0^{2\pi}r u_r(r,\theta)\,\dd\theta
  =\pi(1-\ee^{-r^2}).
\]
代入 (1)，得到
\begin{align*}
  I
  &=\pi\int_0^1 r(1-\ee^{-r^2})\,\dd r\\
  &=\frac{\pi}{2}-\frac{\pi}{2}(1-\ee^{-1})\\
  &=\frac{\pi}{2\ee}.
\end{align*}
故
\[
  \boxed{I=\frac{\pi}{2\ee}}.
\]
\end{xsolution}

\section*{26.3 参考题——第二组参考题}

\subsection*{第二组参考题 1}
\begin{xsolution}
取任意闭球
\[
  \overline{G}\subset\Omega.
\]
以 $u$ 在 $\partial G$ 上的值为边值，用球上的 Poisson 积分公式定义函数 $v$。则 $v$ 在 $G$ 内调和，在 $\overline G$ 上连续，并且
\[
  v=u\qquad\text{于 }\partial G.
\]
调和函数满足平均值公式，所以 $v$ 满足平均值公式；题设又说明 $u$ 满足平均值公式。因此
\[
  w=u-v
\]
也满足平均值公式，且
\[
  w=0\qquad\text{于 }\partial G.
\]

先说明满足平均值公式的连续函数同样满足极值原理。若 $w$ 在 $G$ 内某点 $P$ 取得正的最大值 $M$，则对任意充分小的 $r>0$，
\[
  M=w(P)=\frac{1}{\abs{\partial B_r(P)}}
  \int_{\partial B_r(P)}w\,\dd S.
\]
由于球面上处处有 $w\le M$，上式只能在整个球面上 $w=M$ 时成立。让 $r$ 在允许范围内变化，可知 $P$ 的一个邻域上均有 $w=M$。因此最大值点集既开又闭；由 $G$ 连通且边界上 $w=0$，不可能有 $M>0$。同理，$w$ 也不可能在内部取得负的最小值。于是
\[
  w\equiv0\qquad\text{于 }G.
\]
故 $u=v$ 于 $G$，从而 $u$ 在 $G$ 内调和。

由于 $G$ 可任取，得到
\[
  \Delta u=0\qquad\text{于 }\Omega,
\]
即 $u$ 是调和函数。
\end{xsolution}

\subsection*{第二组参考题 2}
\begin{xsolution}
由边界上的一致收敛，$u_n$ 在 $\partial B_R$ 上一致 Cauchy。对任意 $m,n$，函数
\[
  w_{m,n}=u_m-u_n
\]
在 $B_R$ 上调和并在闭圆盘上连续。由调和函数的极值原理分别用于 $w_{m,n}$ 与 $-w_{m,n}$，
\[
  \sup_{\overline{B_R}}\abs{u_m-u_n}
  \le\sup_{\partial B_R}\abs{u_m-u_n}.
\]
右端随 $m,n\to\infty$ 趋于零，所以 $u_n$ 在 $\overline{B_R}$ 上一致 Cauchy，从而一致收敛到某个连续函数 $u$。

还需证明 $u$ 调和。任取闭圆盘
\[
  \overline{B_\rho(P)}\subset B_R.
\]
每个 $u_n$ 都满足平均值公式
\[
  u_n(P)=\frac{1}{2\pi}\int_0^{2\pi}u_n(P+\rho\vect{e}_\theta)\,\dd\theta.
\]
由一致收敛可以令 $n\to\infty$ 穿过积分号，得到
\[
  u(P)=\frac{1}{2\pi}\int_0^{2\pi}u(P+\rho\vect{e}_\theta)\,\dd\theta.
\]
因此 $u$ 处处满足平均值公式。由第二组参考题 1，$u$ 是调和函数。
\end{xsolution}

\subsection*{第二组参考题 3}
\begin{xsolution}
固定任意 $M_0\in D$，并定义球面平均值
\[
  F(R)=\frac{1}{4\pi R^2}\oiint_{\partial B_R(M_0)}u\,\dd S,
\]
其中 $B_R(M_0)\subset D$。与例题 26.2.2 的 (26.18) 同样计算可得
\begin{equation}
  F'(R)=\frac{1}{4\pi R^2}\iiint_{B_R(M_0)}\Delta u\,\dd V.
  \tag{1}
\end{equation}
并且由连续性
\[
  \lim_{R\to0^+}F(R)=u(M_0).
\]

若 $\Delta u\ge0$ 于 $D$，则由 (1) 有 $F'(R)\ge0$，所以
\[
  F(R)\ge F(0)=u(M_0),
\]
即
\[
  u(M_0)\le\frac{1}{4\pi R^2}
  \oiint_{\partial B_R(M_0)}u\,\dd S.
\]

反过来，假设上述不等式对所有闭球都成立。若存在一点 $M_0$ 使
\[
  \Delta u(M_0)<0,
\]
由连续性可取 $R_0>0$，使 $B_{R_0}(M_0)\subset D$ 且球内 $\Delta u<0$。于是对 $0<R<R_0$，由 (1)
\[
  F'(R)<0.
\]
故 $F(R)<F(0)=u(M_0)$，与题设平均值不等式矛盾。因此不存在这样的点，必有
\[
  \Delta u\ge0\quad\text{于 }D.
\]
充分性与必要性均得证。
\end{xsolution}

\subsection*{第二组参考题 4}
\begin{xsolution}
固定 $M=(x,y,z)\in\Omega$，积分变量记为
\[
  M'=(\xi,\eta,\zeta),
  \qquad
  \vect{r}=M'-M,
  \qquad
  r=\abs{\vect{r}}.
\]
函数 $1/r$ 在 $M'\ne M$ 时调和。取足够小的 $\varepsilon>0$，使闭球 $\overline{B_\varepsilon(M)}\subset\Omega$，并在穿孔区域
\[
  \Omega_\varepsilon=\Omega\setminus\overline{B_\varepsilon(M)}
\]
上对 $u$ 与 $1/r$ 使用第二 Green 恒等式。由于两者在 $\Omega_\varepsilon$ 上都调和，
\[
  0=\oiint_{\partial\Omega_\varepsilon}\left[
  \frac1r\frac{\partial u}{\partial n}
  -u\frac{\partial}{\partial n}\left(\frac1r\right)
  \right]\dd S.
\]
边界由外边界 $S$ 和内球面 $S_\varepsilon=\partial B_\varepsilon(M)$ 组成。

在外边界 $S$ 上，
\[
  \frac{\partial}{\partial n}\left(\frac1r\right)
  =\nabla_{M'}\left(\frac1r\right)\cdot\vect{n}
  =-\frac{\vect{r}\cdot\vect{n}}{r^3}
  =-\frac{\cos(\vect{r},\vect{n})}{r^2}.
\]
因此外边界贡献为
\[
  \oiint_S\left[
  \frac1r\frac{\partial u}{\partial n}
  +u\frac{\cos(\vect{r},\vect{n})}{r^2}
  \right]\dd S.
\]

注意在内球面上，作为穿孔区域 $\Omega_\varepsilon$ 的外法向量是径向向内的 $-\vect{e}_r$，故
\[
  \frac{\partial}{\partial n}\left(\frac1r\right)=\frac1{\varepsilon^2}.
\]
于是内球面贡献为
\[
  \oiint_{S_\varepsilon}\left[
  \frac1\varepsilon\frac{\partial u}{\partial n}
  -\frac{u}{\varepsilon^2}
  \right]\dd S.
\]
第一项的绝对值为 $\bigO(\varepsilon)$；第二项由 $u$ 的连续性满足
\[
  -\frac1{\varepsilon^2}\oiint_{S_\varepsilon}u\,\dd S
  \longrightarrow-4\pi u(M).
\]
令 $\varepsilon\to0^+$，得到
\[
  \oiint_S\left[
  \frac1r\frac{\partial u}{\partial n}
  +u\frac{\cos(\vect{r},\vect{n})}{r^2}
  \right]\dd S
  =4\pi u(M).
\]
故
\[
  \boxed{
  u(x,y,z)=\frac1{4\pi}\oiint_S\left[
  u(\xi,\eta,\zeta)\frac{\cos(\vect{r},\vect{n})}{r^2}
  +\frac1r\frac{\partial u(\xi,\eta,\zeta)}{\partial n}
  \right]\dd S
  }.
\]
\end{xsolution}

\subsection*{第二组参考题 5}
\begin{xsolution}
先把圆心平移到原点。设 $u$ 在 $B_R(0)$ 上非负调和，点 $(x,y)$ 的极坐标为 $(r,\theta)$。若 $u$ 未必连续到半径 $R$ 的边界，则先固定任意
\[
  r<\rho<R.
\]
在闭圆盘 $\overline{B_\rho(0)}$ 上应用 Poisson 公式：
\[
  u(r,\theta)=\frac1{2\pi}\int_0^{2\pi}
  u(\rho,\varphi)
  \frac{\rho^2-r^2}{\rho^2-2\rho r\cos(\theta-\varphi)+r^2}\,\dd\varphi.
\]
由于
\[
  (\rho-r)^2
  \le \rho^2-2\rho r\cos(\theta-\varphi)+r^2
  \le(\rho+r)^2,
\]
Poisson 核满足
\[
  \frac{\rho-r}{\rho+r}
  \le
  \frac{\rho^2-r^2}{\rho^2-2\rho r\cos(\theta-\varphi)+r^2}
  \le
  \frac{\rho+r}{\rho-r}.
\]
又 $u(\rho,\varphi)\ge0$，所以积分后
\[
  \frac{\rho-r}{\rho+r}\cdot
  \frac1{2\pi}\int_0^{2\pi}u(\rho,\varphi)\,\dd\varphi
  \le u(r,\theta)
  \le
  \frac{\rho+r}{\rho-r}\cdot
  \frac1{2\pi}\int_0^{2\pi}u(\rho,\varphi)\,\dd\varphi.
\]
由中心的平均值公式，
\[
  \frac1{2\pi}\int_0^{2\pi}u(\rho,\varphi)\,\dd\varphi=u(0,0).
\]
因此
\[
  \frac{\rho-r}{\rho+r}u(0,0)
  \le u(x,y)\le
  \frac{\rho+r}{\rho-r}u(0,0).
\]
令 $\rho\to R^-$，得
\[
  \frac{R-r}{R+r}u(0,0)
  \le u(x,y)\le
  \frac{R+r}{R-r}u(0,0).
\]
平移回圆心 $(x_0,y_0)$ 即为题设不等式。
\end{xsolution}

\subsection*{第二组参考题 6}
\begin{xsolution}
设 $u$ 是全平面上的有界调和函数。取 $M>0$ 使
\[
  \abs{u(x,y)}\le M
\]
对所有 $(x,y)$ 成立，并令
\[
  v=M+u\ge0.
\]
则 $v$ 仍是全平面上的调和函数。

固定任意点 $P$，记它到原点的距离为 $r$。对任意 $R>r$，在圆盘 $B_R(0)$ 上应用第二组参考题 5 的不等式，得到
\[
  \frac{R-r}{R+r}v(0)
  \le v(P)\le
  \frac{R+r}{R-r}v(0).
\]
令 $R\to\infty$，两侧系数都趋于 $1$，故
\[
  v(P)=v(0).
\]
由于 $P$ 任意，$v$ 为常数，从而 $u=v-M$ 也是常数。
\end{xsolution}
